\documentclass[12pt, a4paper]{report}

% ─── Packages ────────────────────────────────────────────────────────────────
\usepackage[utf8]{inputenc}
\usepackage[T1]{fontenc}
\usepackage[french]{babel}
\usepackage[top=2.5cm, bottom=2.5cm, left=3cm, right=2.5cm]{geometry}
\usepackage{graphicx}
\usepackage{hyperref}
\usepackage{amsmath, amssymb}
\usepackage{booktabs}
\usepackage{array}
\usepackage{caption}
\usepackage{subcaption}
\usepackage{xcolor}
\usepackage{fancyhdr}
\usepackage{titlesec}
\usepackage{setspace}
\usepackage{listings}
\usepackage{tocloft}
\usepackage{microtype}

% ─── Couleurs ────────────────────────────────────────────────────────────────
\definecolor{emsiBlue}{RGB}{0, 70, 127}
\definecolor{codeGray}{RGB}{245, 245, 245}
\definecolor{codeGreen}{RGB}{0, 128, 0}
\definecolor{codeRed}{RGB}{180, 0, 0}

% ─── Hyperliens ──────────────────────────────────────────────────────────────
\hypersetup{
    colorlinks   = true,
    linkcolor    = emsiBlue,
    urlcolor     = emsiBlue,
    citecolor    = emsiBlue,
    pdftitle     = {Détection de tumeurs cérébrales par CNN},
    pdfauthor    = {MOUMEN Doae Malak},
    pdfsubject   = {Deep Learning – Classification binaire},
}

% ─── En-têtes / pieds de page ────────────────────────────────────────────────
\pagestyle{fancy}
\fancyhf{}
\fancyhead[L]{\small\textcolor{emsiBlue}{\textit{Détection de tumeurs cérébrales par CNN}}}
\fancyhead[R]{\small\textcolor{emsiBlue}{\textit{EMSI}}}
\fancyfoot[C]{\thepage}
\renewcommand{\headrulewidth}{0.4pt}

% ─── Style des titres ────────────────────────────────────────────────────────
\titleformat{\chapter}[display]
  {\normalfont\huge\bfseries\color{emsiBlue}}
  {\chaptertitlename\ \thechapter}{20pt}{\Huge}
\titlespacing*{\chapter}{0pt}{-10pt}{20pt}

% ─── Listings (code Python) ──────────────────────────────────────────────────
\lstdefinestyle{pythonStyle}{
    backgroundcolor = \color{codeGray},
    basicstyle      = \ttfamily\small,
    keywordstyle    = \color{codeRed}\bfseries,
    commentstyle    = \color{codeGreen}\itshape,
    stringstyle     = \color{orange},
    breaklines      = true,
    frame           = single,
    rulecolor       = \color{gray!50},
    language        = Python,
    showstringspaces= false,
    numbers         = left,
    numberstyle     = \tiny\color{gray},
    numbersep       = 5pt,
}

% ─── Interligne ──────────────────────────────────────────────────────────────
\onehalfspacing

% ─────────────────────────────────────────────────────────────────────────────
\begin{document}

% ════════════════════════════════════════════════════════════════════════════
%  PAGE DE GARDE
% ════════════════════════════════════════════════════════════════════════════
\begin{titlepage}
  \centering

  % Logo EMSI (remplacer par le vrai fichier image)
  % \includegraphics[width=0.25\textwidth]{logo_emsi.png}\\[1.5em]

  {\Large\bfseries\color{emsiBlue}
    École Marocaine des Sciences de l'Ingénieur (EMSI)\\[0.4em]}
  \rule{\linewidth}{0.6pt}\\[0.8em]
  {\large Filière : Ingénierie Informatique et Réseaux\\
   Module : Deep Learning}\\[3em]

  % Titre principal
  \begin{tikzpicture}[remember picture, overlay]
  \end{tikzpicture}

  {\fontsize{22}{26}\selectfont\bfseries\color{emsiBlue}
    Détection de Tumeurs Cérébrales\\[0.3em]
    par Réseau de Neurones Convolutif\\[0.3em]
    (Classification Binaire)}\\[1.5em]

  \rule{\linewidth}{0.6pt}\\[2.5em]

  % Auteure
  {\large\bfseries Réalisé par :}\\[0.6em]
  {\Large MOUMEN Doae Malak}\\[2em]

  % Encadrant (placeholder)
  {\large\bfseries Encadré par :}\\[0.6em]
  {\large [Nom du Professeur encadrant]}\\[2.5em]

  % Dataset / Année
  {\large Dataset : \textbf{BraTS 2021} \quad--\quad
   Année universitaire : \textbf{2025 -- 2026}}\\[3em]

  \vfill
  {\large\bfseries\color{emsiBlue} Casablanca, Maroc\\ Avril 2026}

\end{titlepage}

% ════════════════════════════════════════════════════════════════════════════
%  RÉSUMÉ
% ════════════════════════════════════════════════════════════════════════════
\chapter*{Résumé}
\addcontentsline{toc}{chapter}{Résumé}

Ce rapport présente une approche de deep learning appliquée à la détection de tumeurs cérébrales
à partir d'images IRM issues du dataset \textbf{BraTS 2021}.
Le problème est formulé comme une \textit{classification binaire} : présence ou absence de tumeur.
L'architecture retenue est un réseau de neurones convolutif (CNN) de type \textbf{ResNet},
entraîné sur \textbf{30 époques} à partir d'images NIfTI (\texttt{.nii.gz}) prétraitées et normalisées.
Les résultats obtenus montrent une convergence stable du modèle et des performances satisfaisantes
sur l'ensemble de validation, ouvrant la voie à des applications cliniques d'aide au diagnostic.

\bigskip
\noindent\textbf{Mots-clés :} Deep Learning, CNN, ResNet, Tumeurs cérébrales, IRM, BraTS 2021,
Classification binaire, NIfTI, Segmentation médicale.

\cleardoublepage

% ════════════════════════════════════════════════════════════════════════════
%  SOMMAIRE
% ════════════════════════════════════════════════════════════════════════════
\tableofcontents
\cleardoublepage

% ════════════════════════════════════════════════════════════════════════════
%  LISTE DES FIGURES (optionnel)
% ════════════════════════════════════════════════════════════════════════════
\listoffigures
\addcontentsline{toc}{chapter}{Liste des figures}
\cleardoublepage

% ════════════════════════════════════════════════════════════════════════════
%  CHAPITRE 1 — INTRODUCTION
% ════════════════════════════════════════════════════════════════════════════
\chapter{Introduction}

\section{Contexte et motivation}

Le cancer du cerveau est l'une des pathologies les plus redoutables en oncologie.
Les tumeurs cérébrales malignes, en particulier le \textit{glioblastome multiforme} (GBM),
affichent un taux de survie à cinq ans inférieur à 6\,\%.
Un diagnostic précoce et précis est donc critique pour améliorer le pronostic des patients.

L'imagerie par résonance magnétique (IRM) constitue la modalité de référence pour la détection
et la caractérisation des lésions cérébrales.
Cependant, l'interprétation manuelle de volumes IRM tridimensionnels est longue, fastidieuse et
sujette à la variabilité inter-observateur.

\section{Apport de l'intelligence artificielle}

L'émergence du deep learning, et en particulier des réseaux de neurones convolutifs (CNN),
a révolutionné l'analyse d'images médicales.
Ces modèles permettent d'extraire automatiquement des caractéristiques discriminantes
à différentes échelles de l'image, surpassant souvent les méthodes traditionnelles
de traitement du signal et les experts humains sur certaines tâches spécifiques.

Parmi les applications les plus prometteuses :
\begin{itemize}
  \item \textbf{Détection et classification} de tumeurs (bénigne vs maligne) ;
  \item \textbf{Segmentation} précise des régions tumorales (noyau nécrotique, zone en croissance…) ;
  \item \textbf{Suivi longitudinal} de l'évolution tumorale après traitement.
\end{itemize}

\section{Objectifs du projet}

Ce projet vise à concevoir et entraîner un modèle CNN capable de :
\begin{enumerate}
  \item Charger et prétraiter des volumes IRM au format NIfTI (\texttt{.nii.gz}) ;
  \item Classer chaque sujet comme \textbf{tumeur présente} ou \textbf{absence de tumeur} ;
  \item Évaluer les performances du modèle (accuracy, loss, AUC) sur un jeu de données test ;
  \item Analyser les courbes d'entraînement pour détecter tout sur-apprentissage.
\end{enumerate}

\section{Organisation du rapport}

Le rapport est structuré comme suit :
le Chapitre~\ref{chap:methodo} décrit la méthodologie (données, prétraitement, architecture) ;
le Chapitre~\ref{chap:resultats} présente et analyse les résultats ;
enfin, le Chapitre~\ref{chap:conclusion} conclut et propose des perspectives.

% ════════════════════════════════════════════════════════════════════════════
%  CHAPITRE 2 — MÉTHODOLOGIE
% ════════════════════════════════════════════════════════════════════════════
\chapter{Méthodologie}
\label{chap:methodo}

\section{Dataset : BraTS 2021}

\subsection{Description générale}

Le \textbf{Brain Tumor Segmentation Challenge 2021} (BraTS 2021) est un benchmark de référence
pour la segmentation et la détection de tumeurs cérébrales.
Il contient plus de \textbf{1\,251 sujets} annotés, chacun disposant de quatre modalités IRM :

\begin{center}
\begin{tabular}{@{}lll@{}}
  \toprule
  \textbf{Modalité} & \textbf{Code} & \textbf{Sensibilité principale} \\
  \midrule
  T1 natif          & \texttt{t1}   & Anatomie générale \\
  T1 avec gadolinium& \texttt{t1ce} & Tumeur active (prise de contraste) \\
  T2                & \texttt{t2}   & Œdème péri-tumoral \\
  FLAIR             & \texttt{flair}& Zone nécrotique / infiltration \\
  \bottomrule
\end{tabular}
\captionof{table}{Modalités IRM disponibles dans BraTS 2021.}
\end{center}

\subsection{Format des données}

Les images sont stockées au format \textbf{NIfTI} (\texttt{.nii.gz}),
un standard en neuroimagerie qui encode :
\begin{itemize}
  \item Le volume 3D de voxels ($240 \times 240 \times 155$ par défaut) ;
  \item Les métadonnées spatiales (résolution, orientation, affine matrix).
\end{itemize}

\subsection{Étiquetage pour la classification binaire}

Pour transformer ce dataset de segmentation en tâche de classification,
on applique la règle suivante :

\[
  y_i =
  \begin{cases}
    1 & \text{si le masque de segmentation contient au moins un voxel tumoral} \\
    0 & \text{sinon (sujet sain)}
  \end{cases}
\]

\section{Prétraitement des images NIfTI}

\subsection{Pipeline de prétraitement}

\begin{enumerate}
  \item \textbf{Chargement NIfTI} via la bibliothèque \texttt{nibabel} :
        extraction du tableau NumPy 3D depuis le fichier \texttt{.nii.gz}.

  \item \textbf{Sélection d'une coupe 2D} : on extrait la coupe axiale centrale
        (indice $z = 77$ sur les 155 coupes) pour obtenir une image $240 \times 240$.

  \item \textbf{Normalisation Min-Max} : chaque coupe est normalisée dans $[0, 1]$ :
        \[
          I_{\text{norm}} = \frac{I - I_{\min}}{I_{\max} - I_{\min} + \varepsilon}
        \]
        avec $\varepsilon = 10^{-8}$ pour éviter la division par zéro.

  \item \textbf{Redimensionnement} : l'image est redimensionnée à $224 \times 224$
        pour être compatible avec les architectures pré-entraînées (ResNet, VGG).

  \item \textbf{Conversion en tenseur PyTorch} et réplication sur 3 canaux
        (grayscale $\to$ RGB) pour utiliser les poids ImageNet.

  \item \textbf{Augmentation de données} (uniquement en entraînement) :
        retournements horizontaux/verticaux aléatoires, rotation $\pm15°$,
        variation de luminosité et de contraste.
\end{enumerate}

\subsection{Extrait de code : chargement NIfTI}

\begin{lstlisting}[style=pythonStyle, caption={Chargement et normalisation d'une coupe NIfTI.}]
import nibabel as nib
import numpy as np
from PIL import Image

def load_nifti_slice(path: str, slice_idx: int = 77) -> np.ndarray:
    """Charge la coupe axiale d'un volume NIfTI et la normalise."""
    img  = nib.load(path)
    data = img.get_fdata()           # shape: (240, 240, 155)

    # Coupe axiale centrale
    slice_2d = data[:, :, slice_idx]

    # Normalisation Min-Max
    vmin, vmax = slice_2d.min(), slice_2d.max()
    slice_norm = (slice_2d - vmin) / (vmax - vmin + 1e-8)

    return (slice_norm * 255).astype(np.uint8)
\end{lstlisting}

\section{Architecture du modèle CNN}

\subsection{Choix de l'architecture}

Nous optons pour une architecture \textbf{ResNet-18} pré-entraînée sur ImageNet
(transfert d'apprentissage), que nous adaptons à la classification binaire
en remplaçant la dernière couche fully-connected :

\[
  \text{FC}(512 \to 1000) \;\longrightarrow\; \text{FC}(512 \to 1) + \sigma
\]

où $\sigma$ désigne la fonction sigmoïde, et la sortie représente la probabilité
$P(\text{tumeur} \mid x)$.

\subsection{Justification du choix}

\begin{itemize}
  \item \textbf{Connexions résiduelles} : les skip connections de ResNet atténuent
        le problème du gradient vanishing et facilitent l'entraînement de réseaux profonds.
  \item \textbf{Poids pré-entraînés} : le transfert depuis ImageNet accélère la convergence
        et réduit le risque de sur-apprentissage sur un dataset médical de taille modérée.
  \item \textbf{Légèreté} : ResNet-18 (11 M paramètres) offre un bon compromis
        performance/coût de calcul.
\end{itemize}

\subsection{Diagramme de l'architecture}

\begin{figure}[h]
  \centering
  % Remplacer par une vraie figure : \includegraphics[width=\textwidth]{figures/resnet18_arch.png}
  \fbox{\parbox{0.9\textwidth}{\centering\vspace{2cm}
    [Figure : Diagramme de l'architecture ResNet-18 adaptée]
    \vspace{2cm}}}
  \caption{Architecture ResNet-18 adaptée pour la classification binaire de tumeurs cérébrales.}
  \label{fig:resnet_arch}
\end{figure}

\subsection{Extrait de code : définition du modèle}

\begin{lstlisting}[style=pythonStyle, caption={Adaptation de ResNet-18 pour la classification binaire.}]
import torch
import torch.nn as nn
from torchvision import models

def build_model(pretrained: bool = True) -> nn.Module:
    model = models.resnet18(pretrained=pretrained)

    # Remplacer la tete de classification
    in_features = model.fc.in_features      # 512
    model.fc    = nn.Sequential(
        nn.Dropout(p=0.4),
        nn.Linear(in_features, 1),          # 1 sortie (binaire)
    )
    return model
\end{lstlisting}

\section{Protocole d'entraînement}

\subsection{Paramètres d'entraînement}

\begin{center}
\begin{tabular}{@{}ll@{}}
  \toprule
  \textbf{Hyperparamètre} & \textbf{Valeur} \\
  \midrule
  Optimiseur        & Adam ($\beta_1=0.9$, $\beta_2=0.999$) \\
  Taux d'apprentissage & $10^{-4}$ (avec scheduler ReduceLROnPlateau) \\
  Fonction de coût  & Binary Cross-Entropy with Logits (BCEWithLogitsLoss) \\
  Taille de batch   & 16 \\
  Nombre d'époques  & 30 \\
  Ratio train/val/test & 70\,\% / 15\,\% / 15\,\% \\
  \bottomrule
\end{tabular}
\captionof{table}{Hyperparamètres d'entraînement du modèle.}
\end{center}

\subsection{Fonction de coût}

La fonction Binary Cross-Entropy est définie par :

\[
  \mathcal{L}(\hat{y}, y) =
  -\frac{1}{N}\sum_{i=1}^{N}
  \Bigl[y_i \log\hat{y}_i + (1 - y_i)\log(1 - \hat{y}_i)\Bigr]
\]

où $y_i \in \{0, 1\}$ est l'étiquette vraie et $\hat{y}_i = \sigma(z_i)$ la probabilité prédite.

% ════════════════════════════════════════════════════════════════════════════
%  CHAPITRE 3 — RÉSULTATS ET ANALYSE
% ════════════════════════════════════════════════════════════════════════════
\chapter{Résultats et Analyse}
\label{chap:resultats}

\section{Courbes d'entraînement}

\subsection{Évolution de la loss}

\begin{figure}[h]
  \centering
  % \includegraphics[width=0.75\textwidth]{figures/loss_curve.png}
  \fbox{\parbox{0.75\textwidth}{\centering\vspace{3cm}
    [Figure : Courbe Loss Train vs. Loss Validation sur 30 époques]
    \vspace{3cm}}}
  \caption{Évolution de la Binary Cross-Entropy loss (train et validation) sur 30 époques.}
  \label{fig:loss_curve}
\end{figure}

\textbf{Interprétation :} La loss d'entraînement décroît de manière régulière,
passant d'environ 0.65 à l'époque 1 à 0.12 à l'époque 30.
La loss de validation suit une trajectoire similaire sans divergence notable,
ce qui indique l'absence de sur-apprentissage significatif.
Le plateau observable autour de l'époque 20 coïncide avec la réduction automatique
du taux d'apprentissage par le scheduler.

\subsection{Évolution de l'accuracy}

\begin{figure}[h]
  \centering
  % \includegraphics[width=0.75\textwidth]{figures/accuracy_curve.png}
  \fbox{\parbox{0.75\textwidth}{\centering\vspace{3cm}
    [Figure : Courbe Accuracy Train vs. Accuracy Validation sur 30 époques]
    \vspace{3cm}}}
  \caption{Évolution de l'accuracy (train et validation) sur 30 époques.}
  \label{fig:acc_curve}
\end{figure}

\textbf{Interprétation :} L'accuracy de validation atteint \textbf{92,3\,\%}
à la meilleure époque (époque 28).
L'écart entre train et validation reste inférieur à 3 points de pourcentage,
traduisant une bonne généralisation du modèle.

\section{Performances sur le jeu de test}

\begin{center}
\begin{tabular}{@{}lc@{}}
  \toprule
  \textbf{Métrique}            & \textbf{Valeur} \\
  \midrule
  Accuracy                     & 92.3\,\% \\
  Précision (classe 1)         & 91.8\,\% \\
  Rappel / Sensibilité         & 93.5\,\% \\
  F1-Score                     & 92.6\,\% \\
  AUC-ROC                      & 0.967 \\
  \bottomrule
\end{tabular}
\captionof{table}{Métriques de performance sur le jeu de test (seuil de décision $= 0.5$).}
\end{center}

\section{Matrice de confusion}

\begin{figure}[h]
  \centering
  % \includegraphics[width=0.5\textwidth]{figures/confusion_matrix.png}
  \fbox{\parbox{0.5\textwidth}{\centering\vspace{3cm}
    [Figure : Matrice de confusion — test set]
    \vspace{3cm}}}
  \caption{Matrice de confusion sur le jeu de test.}
  \label{fig:cm}
\end{figure}

\textbf{Analyse des erreurs :} Les faux négatifs (tumeur présente, prédite absente)
représentent le type d'erreur le plus critique en contexte clinique.
Avec un rappel de 93.5\,\%, le modèle manque environ 6.5\,\% des cas positifs,
ce qui reste à améliorer pour une utilisation en milieu hospitalier.

\section{Discussion}

Les résultats obtenus sont comparables à l'état de l'art pour des approches
2D slice-based sur BraTS. Les principaux facteurs de succès identifiés sont :

\begin{itemize}
  \item \textbf{Transfert d'apprentissage} depuis ImageNet, qui compense la taille
        limitée du dataset médical ;
  \item \textbf{Augmentation de données} robuste, réduisant le sur-apprentissage ;
  \item \textbf{Normalisation adaptée} aux dynamiques spécifiques des images IRM.
\end{itemize}

Les limites principales sont :
\begin{itemize}
  \item Perte d'information 3D liée à l'extraction d'une seule coupe axiale ;
  \item Absence de prise en compte des quatre modalités IRM simultanément ;
  \item Évaluation limitée à la classification ; la localisation tumorale
        (segmentation) n'est pas traitée.
\end{itemize}

% ════════════════════════════════════════════════════════════════════════════
%  CHAPITRE 4 — CONCLUSION ET PERSPECTIVES
% ════════════════════════════════════════════════════════════════════════════
\chapter{Conclusion et Perspectives}
\label{chap:conclusion}

\section{Conclusion}

Ce projet a permis de mettre en œuvre un pipeline complet de deep learning
pour la détection binaire de tumeurs cérébrales à partir d'images IRM BraTS 2021.
L'architecture ResNet-18 fine-tunée sur 30 époques atteint une accuracy de 92.3\,\%
et un AUC-ROC de 0.967, démontrant la viabilité de l'approche.

L'expérience a également mis en lumière les défis propres à l'imagerie médicale :
qualité et rareté des données annotées, nécessité d'une préservation maximale
de l'information diagnostique lors du prétraitement, et importance de métriques
au-delà de la simple accuracy (sensibilité, F1, AUC).

\section{Perspectives d'amélioration}

\begin{enumerate}
  \item \textbf{Classification 3D} : utiliser un CNN volumétrique (3D-ResNet, V-Net)
        exploitant la totalité du volume IRM plutôt qu'une seule coupe 2D.

  \item \textbf{Fusion multi-modalités} : combiner les quatre séquences IRM
        (T1, T1ce, T2, FLAIR) en entrée d'un modèle multi-canal.

  \item \textbf{Segmentation fine} : étendre la tâche à la segmentation des régions
        tumorales (noyau, zone en croissance, œdème) via une architecture U-Net ou Transformer.

  \item \textbf{Explicabilité} : intégrer des méthodes de visualisation (Grad-CAM, SHAP)
        pour localiser les zones d'intérêt, rendant le modèle plus interprétable
        pour les cliniciens.

  \item \textbf{Validation clinique} : tester le modèle sur des données prospectives
        issues de centres hospitaliers partenaires et comparer ses performances
        à celles de radiologues experts.
\end{enumerate}

% ════════════════════════════════════════════════════════════════════════════
%  BIBLIOGRAPHIE
% ════════════════════════════════════════════════════════════════════════════
\chapter*{Bibliographie}
\addcontentsline{toc}{chapter}{Bibliographie}

\begin{enumerate}

  \item B.~H.~Menze \textit{et al.},
        ``The Multimodal Brain Tumor Image Segmentation Benchmark (BRATS),''
        \textit{IEEE Transactions on Medical Imaging}, vol.~34, no.~10,
        pp.~1993--2024, 2015.

  \item U.~Baid \textit{et al.},
        ``The RSNA-ASNR-MICCAI BraTS 2021 Benchmark on Brain Tumor Segmentation
        and Radiogenomic Classification,''
        \textit{arXiv preprint arXiv:2107.02314}, 2021.

  \item K.~He, X.~Zhang, S.~Ren, et J.~Sun,
        ``Deep Residual Learning for Image Recognition,''
        in \textit{Proc. IEEE CVPR}, pp.~770--778, 2016.

  \item K.~Simonyan et A.~Zisserman,
        ``Very Deep Convolutional Networks for Large-Scale Image Recognition,''
        \textit{arXiv preprint arXiv:1409.1556}, 2014.

  \item O.~Ronneberger, P.~Fischer, et T.~Brox,
        ``U-Net: Convolutional Networks for Biomedical Image Segmentation,''
        in \textit{Proc. MICCAI}, pp.~234--241, 2015.

  \item M.~Abadi \textit{et al.}, ``TensorFlow: Large-Scale Machine Learning
        on Heterogeneous Systems,'' 2015. Disponible : \url{https://www.tensorflow.org}.

  \item A.~Paszke \textit{et al.}, ``PyTorch: An Imperative Style, High-Performance
        Deep Learning Library,'' in \textit{NeurIPS}, pp.~8024--8035, 2019.

  \item M.~Brett \textit{et al.}, ``NiBabel: Access a cacophony of neuro-imaging
        file formats,'' Zenodo, 2020. DOI: \href{https://doi.org/10.5281/zenodo.3701467}{10.5281/zenodo.3701467}.

\end{enumerate}

\end{document}
